% Include-only proof fragment for the unchanged abstract K218 source.
% This does NOT certify the rejected shortened-scanner binary table.
% Uses amsmath, amssymb and amsthm; theorem/lemma environments may be supplied
% by the main article. All code constructors are numerical, not syntax tests.
\section{Why the changed startup seed is sufficient}
\label{sec:seed217}

This section proves a theorem about the \emph{unchanged abstract source
program} and its native proof calculus. A proposed shortened-scanner
217-state binary table reaches a tape representing the following clean
root seed during its blank run:
\[
 (N,T,W,A,D,H,X,Y,L)=(4,3,0,0,3,0,0,0,0).
\]
Here $N$ is the next candidate, $T$ the current formula, $W$ the theorem
stack, $A$ the selector, $D$ the additive parameter, $H$ the cached formula,
$X,Y$ the scratch registers, and $L$ the current proof-list pointer.
The zero selector and scratches meet the inherited clean-root interface.
However, that binary table subsequently fails the source simulation
obligation. Consequently the theorem below does \emph{not} certify that
217-state table; it isolates a startup issue which would otherwise have
remained a separate proof obligation.
The nonzero current formula is harmless for a precise reason: native code
$3$ is a theorem. No assumption of consistency is used to prove this.

Write $\mathcal C$ for the exact C222 numerical calculus inherited by K218.
It uses
\[
 P(a,b)=\frac{(a+b)(a+b+1)}2+b,\qquad
 I(a,b)=4P(a,b),\qquad R(a,b)=4P(a,b)+1,
\]
\[
 Q(i,a)=4P(i,a)+2,\qquad E(a,b)=4P(a,b)+3.
\]
We write $a\to b=I(a,b)$, $Q_i a=Q(i,a)$,
$\neg a=I(a,1)$ and $\exists_i a=\neg Q_i\neg a$.
All parameters range over natural codes, including any values which would
be considered junk by a separate syntax. Truth is code zero, falsum is
code one, and the alias $I(0,0)=0$ remains literal.

\subsection{A concrete proof that code three is a theorem}

\begin{lemma}[Native equality reflexivity]
For every natural $x$, $\mathcal C\vdash E(x,x)$. In particular,
$\mathcal C\vdash 3$.
\end{lemma}
\begin{proof}
The inherited unrestricted $\rho_8$ supplies the derived K and S schemata,
and consequently the propositional deduction transformation. For closed
theorems $a\to b$, the retained $\rho_2$ row and theorem-only generalization
give $Q_i a\to Q_i b$. Applying this to the propositional contraposition
of $a\to b$, then contraposing once more, gives the closed derived rule
\[
 \frac{\mathcal C\vdash a\to b}
      {\mathcal C\vdash \exists_i a\to\exists_i b}.
\]
Every use of generalization in this argument is on a closed theorem.

Put $r=R(x,x)$, $e=E(x,x)$ and $p=E(r,x)$. The merged equality row at
$(a,b,c)=(x,x,r)$ is exactly
\[
 p\to(p\to Q_r e).
\]
By propositional contraction, $\mathcal C\vdash p\to Q_r e$.
Existential monotonicity gives
$\mathcal C\vdash \exists_r p\to\exists_r Q_r e$.
The retained $\rho_1(r,x)$ is $\exists_r p$, so modus ponens yields
$\exists_r Q_r e$. Finally $\rho_6(r,e)$ is
$\exists_r Q_r e\to e$, proving $e$.
For $x=0$, the arithmetic definition gives $E(0,0)=3$.
\end{proof}

This is the same ordinary argument formalized upstream as
\texttt{G222.refl222}. The accompanying file
\texttt{seed\_reflexivity\_proof.json} specializes it to zero and retains
only the 79 closed nodes which are ancestors of the conclusion. Its
independent forward checker accepts only literal K/S instances, the
retained $\rho_1,\rho_2,\rho_6,G$ instances, modus ponens, and
generalization of a previously closed theorem. K and S are explicitly
identified as inherited derived schemata; this certificate does not add
them as new machine axioms. There is no hypothesis rule in the retained
certificate and no arbitrary object-language substitution inference.

\subsection{Soundness from any theorem-valued clean seed}

For a natural stack code $w$, define
\[
 \operatorname{Good}(w)\quad\Longleftrightarrow\quad
 \forall k\in\mathbb N,\;
 \mathcal C\vdash
 \pi_1\!\left(\operatorname{unP}
       \bigl((\pi_2\circ\operatorname{unP})^k(w)\bigr)\right).
\]
Thus every entry in the infinite zero-padded stack stream is a theorem.
The inverse identities for Cantor pairing prove that pushing a theorem
preserves Good, and that popping a Good stack produces a theorem and a
Good tail. The stack $0$ is Good because unpairing zero gives $(0,0)$
and zero is a native theorem. These facts include the alias $P(0,0)=0$.

At a command boundary, the invariant is
\[
 \mathcal C\vdash T,\qquad
 \mathcal C\vdash H,\qquad
 \operatorname{Good}(W).
\]
It places no restriction on $N,D,L$. Its preservation is immediate from
the literal command semantics, as follows. A successful MP uses the
theorem cache $H=I(T,D+x)$ and the theorem $T$ to obtain $D+x$;
it pushes zero and clears the cache. A failed MP replaces the current
formula and cache by zero. DROP pushes the theorem $T$ and replaces
it by the theorem $H$. NIP first tests whether the theorem $H$ equals
one; if so, it halts with a native refutation. Otherwise it pops a
theorem from the Good stack into the cache. AX pushes the old theorem
$T$, a universal generalization of the theorem cache, and the eight
literal native axiom rows. Its new current formula is the gated bundle
or zero, and its cache is cleared. Thus every nonhalting command
preserves the invariant, and every halting command proves
$\mathcal C\vdash1$.

For the displayed new startup seed, $T=3$ is covered by reflexivity,
$H=0$ is truth, and $W=0$ is Good. Therefore a halt from that seed
implies a native refutation. The same proof is already present as the
generic \texttt{Inv\_step}, \texttt{Inv\_run}, and \texttt{halt\_step}
lemmas in the inherited \texttt{L1\_Sound222.v}; only its base case is
changed here.

\subsection{Completeness above any finite candidate cutoff}

The changed next-candidate value $N=4$ requires a cofinality argument;
merely quoting a theorem for an all-zero seed would not suffice. The
inherited proof compiler supplies, for every native derivation of a code
$X$, a finite command block $B_X$ with the following strong contract:
from \emph{every} retained pair $(t,w)$, execution starting at
$(T,W,D,H)=(t,w,0,0)$ either halts early or ends exactly at
$(X,w,0,0)$. The ordinary leaf, GEN, MP and discard proofs establishing
this contract do not depend on the initial $t$ or $w$ being theorems.

Suppose $\mathcal C\vdash1$, and fix its block $B_1$. Let
\[
 S=[\operatorname{DROP}(0),\operatorname{NIP}(0),
                         \operatorname{NIP}(0)].
\]
For every natural $u$, use the first command
\[
 \operatorname{START}_u=\operatorname{AX}(0,1,1,u),
\]
where the three actual AX fields after the selector are the common
additive parameter, cached-parameter increment, and unrestricted fourth
argument. In full wire order the fields are $[3,1,1,u]$.
Every AX command resets $D=H=0$, regardless of their incoming values.
Consequently
\[
 [\operatorname{START}_u]\mathbin{+\!+}B_1\mathbin{+\!+}S
\]
halts from every retained tuple $(T,W,D,H)$. If $B_1$ halts, there is
nothing further to show. If it returns normally, its current formula
is one and its cache is zero; the three-command suffix first pushes
one, then loads it into the cache, then halts.

Let $r$ be the fixed encoded command-list tail for $B_1\mathbin{+\!+}S$.
The corresponding initial candidate number is
\[
 n_u=P\bigl(3,P(1,P(1,P(u,r)))\bigr).
\]
Cantor pairing dominates either argument, hence $n_u\geq u$. These are
therefore successful candidate numbers above any finite initial cutoff.
The fact that different $u$ values change one axiom emitted by the
initial AX does not affect the block contract; its arbitrary incoming
top and retained stack are already quantified universally.

The literal scheduler first loads candidate $n$ directly, and decodes
every subsequent command from the predecessor of a positive proof-list
pointer. Each unpair remainder is at most its input. Thus whenever
$L>0$ and the command does not halt, the next pointer is at most $L-1$.
Induction on $L$ drains each current proof list in finitely many
commands, without changing the next-candidate register. When $L=0$,
the scheduler loads $N$ and increments $N$. It follows by induction
that every candidate number at or above any finite initial $N$ is
eventually started, unless a halt occurs first. Choose $u\geq N$;
candidate $n_u$ supplies a halt if no earlier one has occurred.

\begin{theorem}[Safe clean seeds]
At a clean source root with $A=X=Y=0$, assume
$\mathcal C\vdash T$, $\mathcal C\vdash H$, and $\operatorname{Good}(W)$.
Allow arbitrary natural $N,D,L$. The source halts from this seed if and
only if $\mathcal C\vdash1$.
\end{theorem}
\begin{proof}
Soundness is the invariant above. Completeness is the strong compiler,
the resetting $\operatorname{START}_u$ family, and the cofinal scheduler
argument. The proof uses finite syntax, natural arithmetic and induction,
so it is formalizable in ZFC. It does not assume $\operatorname{Con}(\mathrm{ZFC})$.
\end{proof}

Applying this theorem to $(4,3,0,0,3,0,0,0,0)$ proves the required source
calibration. The unchanged K218 source-to-calculus correspondence and
inherited native-calculus/ZFC refutation translations then give
\[
 \operatorname{Halts}(\text{source},(4,3,0,0,3,0,0,0,0))
 \quad\Longleftrightarrow\quad
 \text{there is a finite ZFC refutation}.
\]
To transfer this conclusion to a physical Turing machine, both a correct
whole-tape startup and an unbounded source-to-binary simulation are still
required. The shortened-scanner table has a verified finite startup
trace but fails the latter obligation: concrete primitive executions
return to incorrect source addresses. No 217-state machine theorem
follows from this seed result.
